\scanpage{117}
Prenons $X=\Delta_n$, alors le diagramme (1) devient
\[
\diagram{\begin{tikzcd}[column sep=large,row sep=large]
\mathrm R^*\Gamma(\Delta_n,S_k)\simeq S_k\arrow[r]&\mathrm R^*\Gamma(\Delta_n,C^{**})=C_n^{**}\\
&C^{**}(\Delta_n)=C_n^{**}\arrow[u,"\mathrm{id}_{C_n^{**}}"']
\end{tikzcd}}
\]
et la condition énoncée se réduit à celle concernant la flèche horizontale, qui signifie aussi
\[
(3)\qquad\diagram{\begin{tikzcd}
H_C^{**}(\Delta_n)=H^{**}(C_n^{**})\quad\deleted{\unsure{est isom.}}\\
S_k\arrow[u,"\epsilon\ \text{augmentation}"']
\end{tikzcd}}
\]
l'augmentation est un quasi-iso

i.e.
\[
3\mathrm a)\quad\deleted{\unsure{k=}}S_k^i\longrightarrow
\begin{gathered}H^{i0}(C_n^{**})\\\Vert\\
\operatorname{Ker}(C_n^{i,0}\longrightarrow C_n^{i1}).\end{gathered}
\quad\deleted{\unsure{=0}}\text{ est un}
\]
iso pour $i$ grand ($i\geq\unsure{\lambda(\infty)|\lambda(0)}$ disons)
\[
3\mathrm b)\quad\forall p>0,\quad H^{ip}(C_n^{**})=0\text{ pour }i
\]
$i$ grand (dépendant de $p$) soit
\[
H^{ip}(C_n^{**})=0\qquad\text{si}\quad i\geq\deleted{\unsure{\varphi}}\inserted{\lambda}(p)
\]
où on peut supposer $\lambda:\mathbb N^{\unsure{*}}\longrightarrow\mathbb N^{\unsure{*}}$ strictement croissante.

A priori, ce $\lambda=\lambda_n$ dépend du choix de $n$. On suspecte qu'il faut pouvoir prendre $\lambda$ indépendant de $n$ -
\margin{116}
